\section{Theory: what a start can and cannot change}\label{sec:theory}

The propositions of this section are exact statements about the iteration of Section~\ref{sec:m07} and the network of Section~\ref{sec:network}; the relation of Section~\ref{sec:conjecture} is a conjecture with the evidence for it. Together they say where a warm start can act and why the gain it can give is bounded by the freeze index.

\subsection{Invariances}\label{sec:invariance}

\begin{proposition}[global phase]\label{prop:global}
Let $c\in\mathbb R$ and let the start phases be shifted, $\theta'_m=\theta_m+c$ for all $m$. Then for every $k\ge0$ the holograms satisfy $\phi'^{(k)}=\phi^{(k)}+c$ and the fields $V'^{(k)}=e^{ic}V^{(k)}$, so $\relstd_k$, $\eta_k$ and $u_k$ are the same sequences.
\end{proposition}
\begin{proof}
For $k=0$, Eq.~\eqref{eq:start} gives $\phi'^{(0)}=\arg\big(e^{ic}\sum_ma_me^{i(\Delta^m+\theta_m)}\big)=\phi^{(0)}+c$, and Eq.~\eqref{eq:forward} gives $V'^{(0)}_m=N^{-2}\sum_jA_je^{i(\phi^{(0)}_j+c-\Delta_j^m)}=e^{ic}V^{(0)}_m$. Assume the claim for $k-1$. The relative amplitudes $B_{k-1}(m)=|V^{(k-1)}_m|/a_m$ do not change, so the weights of Eq.~\eqref{eq:weights} do not change; the phases of Eq.~\eqref{eq:freeze} become $\psi'_{k,m}=\psi_{k,m}+c$ (and the frozen
phase $\psi'_{\nfix}=\psi_{\nfix}+c$). Eq.~\eqref{eq:update} then gives $\phi'^{(k)}=\phi^{(k)}+c$, and Eq.~\eqref{eq:forward} gives $V'^{(k)}=e^{ic}V^{(k)}$. The intensities are unchanged. \end{proof}

Proposition~\ref{prop:global} is why the loss of Eq.~\eqref{eq:loss}, which reads intensities only, needs no special treatment of the global phase, and why a regression onto converged phases would need one.

\begin{proposition}[permutation equivariance]\label{prop:equivariance}
Let $\pi$ be a permutation of the valid traps and $\pi\mathcal T$ the permuted trap set. The network of Eqs.~\eqref{eq:features}--\eqref{eq:output} satisfies $g_\vartheta(\pi\mathcal T)=\pi\,g_\vartheta(\mathcal T)$, and the start hologram of Eq.~\eqref{eq:start} is the same for $\mathcal T$ and $\pi\mathcal T$.
\end{proposition}
\begin{proof}
The features of Eq.~\eqref{eq:features} and the embedding act on each trap separately, so they commute with $\pi$. In Eq.~\eqref{eq:attention} the queries are permuted with the traps; the logits of query $m$ over the keys are permuted with the keys, and the softmax and the weighted sum of the values are invariant to a permutation of the keys, so the output for trap $m$ is unchanged; the additive mask follows its trap.
The mean $\bar{\mathbf h}$ over the valid traps is invariant to $\pi$, and the readout of Eq.~\eqref{eq:output} acts row by row with $\bar{\mathbf h}$ and $\varrho_m^2$ attached to trap $m$. Hence $\theta_m$ follows trap $m$. The sum of Eq.~\eqref{eq:start} runs over all traps and does not depend on their order. \end{proof}

The network has no positional index of any kind, so the order in which a user lists the traps cannot change the start, and the padded entries (removed by the mask) cannot change it either. The proposition is also the reason for a set network and against a recurrent or convolutional one: an ordering of traps would be an arbitrary input that the network would have to learn to ignore.

\subsection{What the start changes}\label{sec:start-change}

\begin{proposition}[the start enters through the frozen phase]\label{prop:washout}
(a) For $\nfix=1$, every hologram $\phi^{(k)}$, $k\ge1$, is a function of the trap set and of the vector $\psi_1=\arg V^{(0)}$ alone. (b) For general $\nfix$, every $\phi^{(k)}$ with $k\ge\nfix$ is a function of the trap set and of the trajectory $V^{(0)},\dots,V^{(\nfix-1)}$, that is, of the start through $\nfix-1$ applications of Eqs.~\eqref{eq:weights}--\eqref{eq:update}.
\end{proposition}
\begin{proof}
(a) With $\nfix=1$ Eq.~\eqref{eq:freeze} gives $\psi_{k}=\psi_1$ for all $k\ge1$ and $g_1=\mathbf 1$, so $\phi^{(1)}=\arg\sum_ma_me^{i(\psi_1+\Delta^m)}$ depends on $\psi_1$ only. Then $V^{(1)}$ is a function of $\phi^{(1)}$, $g_2$ of $V^{(1)}$ by Eq.~\eqref{eq:weights}, and $\phi^{(2)}=\arg\sum_mg_{2,m}a_me^{i(\psi_1+\Delta^m)}$ of $g_2$ and $\psi_1$; the induction continues. (b) The frozen phase is $\psi_{\nfix}=\arg V^{(\nfix-1)}$ by Eq.~\eqref{eq:freeze}, and the weights $g_k$ for $k\le\nfix$ are functions of $V^{(1)},\dots,V^{(\nfix-1)}$ by Eq.~\eqref{eq:weights}; after the freeze the sequence is generated by the same two rules from these quantities, so the start enters only
through the trajectory up to $V^{(\nfix-1)}$. \end{proof}

Part (a) turns the warm-start problem at $\nfix=1$ into the prediction of $M$ phases and nothing else: the network's $\theta_m$ become, to the extent that the cross terms of Eq.~\eqref{eq:start} are small, the frozen phases of the whole run. Part (b) holds for any deterministic iteration and constrains nothing by itself; its content here is that the start enters only through $V^{(0)}$ while the weights restart at one, so at the engine default $\nfix=12$ the start has to survive eleven applications of the iteration to matter. Whether it does is a property of the dynamics of the iteration, not of the network, and it is measured in Section~\ref{sec:ablation}: from the phases of a converged
run (the best start one could hope to predict) the iteration at $\nfix=12$ still needs \OTwelveMed{} iterations (median over ten geometries) against \ORMed{} from a random start.

\subsection{An accounting bound on the weights}\label{sec:weights-bound}

The weights restart at one whatever the start is, so a warm start does not transfer the correction that a converged run has stored in $g$. How long the weights need to rebuild it is bounded below by an accounting identity. Let $S_k=\max_m\log g_{k,m}-\min_m\log g_{k,m}$ be the log-spread of the weights after iteration $k$, and
$s_k=\log\big(\max_mB_{k-1}(m)/\min_mB_{k-1}(m)\big)$ the allowance of the update at iteration $k\ge2$.

\begin{proposition}[allowance]\label{prop:allowance}
For every run and every $k\ge2$,
\begin{equation}
S_k\ \le\ \sum_{j=2}^{k}s_j .
\label{eq:allowance}
\end{equation}
\end{proposition}
\begin{proof}
By Eq.~\eqref{eq:weights}, $\log g_{k,m}=\log g_{k-1,m}+\log\langle B_{k-1}\rangle-\log B_{k-1}(m)$. The term $\log\langle B_{k-1}\rangle$ is common to all traps and cancels in the spread. The spread of a sum is at most the sum of the spreads, and the spread of $-\log B_{k-1}(m)$ over $m$ is $s_k$. Hence $S_k\le S_{k-1}+s_k$, and $S_1=0$ since $g_1=\mathbf 1$. \end{proof}

If a converged run holds a weight spread $S^\star$, no run that starts from $g=\mathbf 1$ reaches that spread before the iteration $k_{\mathrm{floor}}=\min\{k:\sum_{j=2}^ks_j\ge S^\star\}$. The bound is a statement about the restart; its tightness is measured, not assumed: we compute $S_k$ and the allowance along the runs of Section~\ref{sec:ablation} and compare $k_{\mathrm{floor}}$ with the iterations to 1 percent (Table~\ref{tab:weights}).

\subsection{A relation between the freeze index and the iterations}\label{sec:conjecture}

\begin{conjecture}[iterations follow the freeze]\label{conj:freeze}
For the random start and the network start on dense arrays (at least 40 traps) for which the iteration reaches 1 percent, the number of iterations to 1 percent is the freeze index plus a short tail, $k_1=\nfix+t$ with $t$ of the order of one to four iterations on most geometries, and the start changes $t$ and the probability of reaching 1 percent. The relation is a description of the starts tried here and has counterexamples, listed below.
\end{conjecture}

We do not prove the conjecture; the evidence is of four kinds, all in Section~\ref{sec:ablation}. (i) Before the freeze the spread of the default run stays at several percent: the median trajectory of the random start at $\nfix=12$ is \trRTwelveEleven{} percent at iteration 11 and \trRTwelveTwelve{} percent at 12, and drops to \trRTwelveThirteen{} percent at 13, the first iteration that uses the frozen phase. (ii) The sweep of the freeze index
puts the median $k_1$ of the random start at \swTwelveEx, \swEightEx, \swSixEx, \swFourEx{} and \swThreeEx{} iterations above the freeze index for $\nfix=12,8,6,4,3$; these medians include the runs that reach the cap of the budget (\swThreeNever{} percent of the runs at $\nfix=3$ and \swFourNever{} percent at $\nfix=4$), and the excess of the runs that converge alone was not computed, because the sweep stores medians and not the count of each run, so the tail of one to four is exceeded at $\nfix=3$ and is not verified on converged runs. (iii) The network start at $\nfix=12$ has the same median as the random start, and (iv) the phases of a converged run lower the median at $\nfix=12$ by \OTwelveGain{} iteration only. The counterexamples are in Table~\ref{tab:oracle}: the random start reaches 1 percent before the freeze index at $\nfix=12$ on one of the ten geometries (\texttt{gen/09}), the converged phases do so on three (\texttt{gen/00}, \texttt{gen/07} and \texttt{gen/09}), and at $\nfix=3$ the converged phases lower the median from \ORThreeMed{} to \OThreeMed{} iterations and converge on three geometries where the random start never does. These are ten geometries and one seed. The conditions under which the relation fails are also in the data: a freeze that is too early locks phases that are
incompatible with a uniform solution, and then the run never reaches 1 percent; the share of such runs rises as the freeze index falls (Table~\ref{tab:sweep}). A warm start can therefore act on the second term, on the share of runs that converge, and on the final spread, and could lower the first term only by lowering the index at which freezing is safe.

\subsection{Consequences for the experiment}\label{sec:consequences}

The propositions fix what the experiments must separate. Because the iteration depends on the start only through $\psi_{\nfix}$ (Proposition~\ref{prop:washout}), a network start and a change of $\nfix$ are two different handles on the same quantity, and an experiment that changes both at once cannot attribute the effect: the candidate network was trained and first evaluated at $\nfix=3$ while its baseline ran at the default $\nfix=12$. The ablation of Section~\ref{sec:ablation} therefore crosses the two factors, and the sweep of $\nfix$ with
and without the network tests Proposition~\ref{prop:washout}(a) directly at $\nfix=1$. A fair comparison of a learned start therefore needs the baseline at the same freeze index as the candidate, in addition to the baseline at the published default: the first isolates the contribution of the network, the second measures the practical gain over the method as published. The first evaluation of the candidate compared it only with the second, and the difference between the two comparisons is the result of this paper. Proposition~\ref{prop:allowance} gives a check that costs nothing: the weights cannot be the binding constraint if $k_{\mathrm{floor}}$ is far below the measured iterations.
